\documentclass{article}

\usepackage[final]{neurips_2025}
\bibliographystyle{plainnat}

\usepackage[T1]{fontenc}
\usepackage[utf8]{inputenc}
\usepackage{graphicx}
\usepackage{booktabs}
\usepackage{amsmath, amssymb}
\usepackage{xcolor}
\usepackage[hidelinks, breaklinks=true]{hyperref}
\usepackage{listings}

\title{From Pull Requests to Sandboxed Restoration:\\
       A Logbook of Building an LLM Hardware-Agent Benchmark}

\author{%
  David Andrews\\
  \texttt{dha@xoba.com}\\
  \And
  Saketh Reddy\\
  \texttt{yssaketh@gmail.com}\\
}

\begin{document}
\maketitle

\begin{abstract}
This logbook documents two and a half months of work building an LLM-agent benchmark for
hardware description languages, ending with the \textsc{RTL-Universe} benchmark we report
in our companion paper. The arc covers four prior repositories --- a GitHub-PR data pipeline
(SWE-Universe-Verilog), an early eval harness (RTLBench), a port of NVIDIA's CVDP problem
suite into our infrastructure (CVDP-Harbor), and a spec-to-GDSII benchmark we built and
iterated weekly through three pivots (Spec2GDSBench) --- before the move to project-scale
\emph{restoration} of open-source HDL projects (\textsc{RTL-Universe}). We write the logbook
as a class progress report: what we tried, what worked, and where we pivoted, with concrete
references to commits, weekly reports, and run results that are still in the project's
\texttt{other\_repos/} directory.
\end{abstract}

\section{Goal}

The high-level question we set out on was: \emph{can current LLM agents do real
hardware-engineering work, and how do we measure that?} The literature on hardware LLMs
in early 2026 was almost entirely module-level Verilog generation~\citep{liu2023verilogeval,
lu2024rtllm, thakur2023verigen}, with leading models reporting 90\%+ pass rates on
benchmarks that ask for ``a 4-bit counter from a description.'' That left two open
questions: how would those agents fare on the rest of a real chip-design workflow (build
systems, integration, testbenches, physical design), and is there a benchmark that we
trust to grade them on it?

Over two and a half months we built five repositories, in roughly chronological order:
\textsc{SWE-Universe-Verilog}, \textsc{RTLBench}, \textsc{CVDP-Harbor},
\textsc{Spec2GDSBench}, and \textsc{RTL-Universe}. Each one taught us something that we
then either folded into the next or used as the reason to start from scratch. The rest of
this document is that arc.

\section{March 2--6: SWE-Universe-Verilog (data pipeline)}
\label{sec:sweuv}

\paragraph{What we tried.}
SWE-bench~\citep{jimenez2024swebench} grades agents on real GitHub bugs in real Python
projects. The first thing we tried was the same recipe but for HDL. The
\textsc{SWE-Universe-Verilog} pipeline crawls GitHub for repositories whose primary
language is Verilog or SystemVerilog (or that match HDL topics: \texttt{verilog},
\texttt{fpga}, \texttt{rtl}, \texttt{hdl}, \texttt{asic}), collects all merged pull
requests, and applies a multi-stage filter inspired by SWE-Universe's approach
\citep{sweuniverse}.

\paragraph{Pipeline yield.}
\begin{center}
\footnotesize
\begin{tabular}{lr}
\toprule
Stage & Count \\
\midrule
Repos discovered (Verilog/SV on GitHub)                & 3{,}296 \\
PRs collected (merged)                                  & 173{,}106 \\
After heuristic filter (size, bot, merge)               & 27{,}446 \\
After testbench detection                               & 5{,}282 \\
After patch separation (test + fix files)               & 3{,}070 \\
Quality filter (relaxed)                                & 1{,}241 \\
Quality filter (strict, with issue linkage)             & 401 \\
\bottomrule
\end{tabular}
\end{center}

\paragraph{What we learned.}
Even after aggressive filtering, the yield from the GitHub HDL ecosystem is roughly an
order of magnitude lower than the Python equivalent. 401 strict-filter PRs across
3,296 repositories is enough to start a benchmark, but each one had to be curated by hand
to be sure it isolated a single change with a real before-and-after test. We built
\textsc{rtlbench} (\S\ref{sec:rtlbench}) on top of this data and quickly discovered that
PR-based hardware bugs do not, in practice, isolate cleanly: most are part of a larger
refactor, depend on FPGA-vendor toolchains, or test only a manually-asserted simulation
trace.

This phase informed two decisions in \textsc{RTL-Universe}: we chose to abandon
PR-replay as the framing, and we chose \emph{restoration} (delete-and-rebuild) instead,
because it gives the agent a real test suite (the upstream's own CI) without requiring us
to curate a clean before-and-after diff.

\section{March 2--3: RTLBench (first agent runs)}
\label{sec:rtlbench}

\paragraph{What we tried.}
Using the SWE-Universe-Verilog filtered set, plus a few additional integration projects
(corundum networking, CVA6 RISC-V core), we built a quick eval harness called
\textsc{RTLBench}: clone the upstream repo at the parent commit of a chosen PR, set up the
toolchain in a container, and ask Claude to make the change. The early commits were
exploratory --- ``init,'' ``snap,'' ``first claude run,'' ``claude plan expand not run
yet,'' ``ran full eval'' (commits \texttt{45e61b9}--\texttt{38a68a1}).

\paragraph{What we learned.}
The first runs worked just well enough to show two failure modes that ended up shaping
everything that followed:
(1)~\textbf{Toolchain fragility}: Verilog projects often pin specific Verilator/Yosys
versions, and the agent would happily install a different one and write code that compiled
in the wrong version. We ended up writing per-task Dockerfiles, which is what the next two
projects (CVDP-Harbor and Spec2GDSBench) inherit.
(2)~\textbf{Test-suite ambiguity}: many of the PRs in our filtered set have testbenches
that ``run'' on master but actually only pass when accompanied by a manually-checked
waveform. The agent does not know this. We needed verifiers whose pass criterion is
mechanically checkable. This is part of why \textsc{RTL-Universe} restricts itself to
projects with binary-pass CI.

\section{March 7--13: CVDP-Harbor}
\label{sec:cvdp}

\paragraph{What we tried.}
At this point NVIDIA released the Comprehensive Verilog Design Problems
(CVDP) benchmark~\citep{nvidia2024cvdp}, a 783-problem suite of mostly module-level Verilog
tasks with both non-agentic and agentic variants. Rather than re-invent the
problems, we ported CVDP into \emph{Harbor}, our internal evaluation harness used by the
rest of the projects. The conversion tool is at \texttt{cvdp\_benchmark/} commit
\texttt{c002eb5}, with three more rounds of fixes and reruns over the next week.

\paragraph{Result.}
We ran three Claude variants on the converted suite. Pass rates clustered above 90\%
(\texttt{cvdp\_benchmark} commit \texttt{891039d}), consistent with what NVIDIA themselves
reported. This was when we became confident that \emph{module-level Verilog generation is
saturated for current frontier models} and that the open question is at the project scale.
At this point we stopped working on CVDP and started a benchmark we hoped would not
saturate.

\section{March 23--April 12: Spec2GDSBench (three weekly pivots)}
\label{sec:spec2gds}

\paragraph{Goal.} Given a behavioural specification and a test suite inside a containerised
environment with the full open-source ASIC toolchain (LibreLane + SkyWater Sky130A PDK),
ask the agent to autonomously produce a DRC-clean, LVS-verified GDSII file. The grading
funnel is: functional simulation passes $\to$ Yosys synthesis $\to$ LibreLane place-and-
route $\to$ KLayout/Magic DRC $\to$ Netgen LVS $\to$ GDSII export. Each stage gates the
next.

\paragraph{Week 1 (Mar 24--28): eight isolated tasks.}
The first version had eight tasks: \texttt{gcd\_unit}, \texttt{uart\_tx},
\texttt{aes128\_encrypt}, \texttt{femtorv32\_rv32i}, and four AI accelerator primitives
(\texttt{conv1d\_engine}, \texttt{maxpool2d}, \texttt{requantize\_pipeline},
\texttt{matmul\_4x4\_int8}). 64 runs of Claude Code Sonnet at \$5/run got mean reward 0.964
in normal mode, 0.829 in blind mode (no test suite visible). All 32 normal-mode runs
passed all functional tests; the only score loss was timing closure during physical
design~\citep{spec2gdsbench-w1}. The benchmark was near-saturated. We pivoted.

\paragraph{Week 2 (Mar 31--Apr 6): single SoC, Linux boot.}
We replaced the eight tasks with one task: design an RV64GC RISC-V SoC capable of booting
Linux. The agent receives only the external interface (AXI4 + UART) and a software
contract (108 ISA compliance tests, 10 bare-metal programs, OpenSBI/BusyBox boot via
QEMU's \texttt{virt} machine convention). Claude Sonnet 4.6 across five runs scored a best
of 0.616 (85/108 ISA, 10/10 bare-metal, OpenSBI boot). A 15-hour run reached 99/108 ISA
and showed partial Linux boot during debugging~\citep{spec2gdsbench-w2}. This was a much
harder benchmark, but a single task gave poor signal: an agent either built a working
SoC or scored near zero. We pivoted again.

\paragraph{Week 3 (Apr 7--12): AI accelerator integration ladder.}
We then designed three new tasks at graduated difficulty, all with AXI4 interfaces:
\texttt{systolic\_array} (8$\times$8 INT8 MAC array, AXI4-Lite MMIO + AXI4-Stream),
\texttt{conv2d\_engine} (3$\times$3 convolution with line buffer, AXI4 master),
\texttt{dma\_engine} (memory-to-memory DMA, AXI4-Lite + AXI4 master). All three reference
solutions pass the full ASIC flow through GDSII; only the DMA meets timing at
50\,MHz~\citep{spec2gdsbench-w3}. We were partway through getting agent runs on these
when the next pivot happened.

\paragraph{What we learned across the three weeks.}
\textbf{Specifications are hard.} Every week spent on Spec2GDSBench, we spent some
fraction of the time on whether the spec was over-constraining the agent (forcing a
specific microarchitecture) or under-constraining it (letting the agent solve a different
problem). Real industrial specs are often closer to ``here is a JIRA epic with 30 linked
tickets'' than ``here are the AXI4 signal names.'' We could not write a clean spec for a
realistic design.

\textbf{Build systems matter as much as RTL.} The agent's score on Spec2GDSBench was
dominated by physical design, not RTL design --- timing closure, place-and-route choices,
LibreLane configuration. We wanted RTL-design ability to be the bottleneck.

\textbf{Real projects already have all of this.} The thing that finally moved us was
realising that lowRISC's Ibex, NVIDIA's NVDLA, OpenPiton, and the rest already have
testbenches, build systems, CI configurations, and reference implementations carefully
tuned by their maintainers. We do not need to write a spec for a UART --- lowRISC has a
real one and grades it with their own simple\_system. We just need to delete files.

\section{April 23 -- May 8: RTL-Universe}

\paragraph{What it is.}
\textsc{RTL-Universe} is the result of all the prior decisions: real upstream projects,
real upstream CI, no spec to write, restoration as the framing.
The task list is in our companion paper. The harness builds on
the Harbor runner that started in CVDP-Harbor (\S\ref{sec:cvdp}); the Docker patterns
follow Spec2GDSBench's three-layer image; the verifier conventions are upstream's CI.

\paragraph{The first three weeks.}
Most of late April and early May was task engineering --- we wrote a
\texttt{harbor-task-gen} skill (commits \texttt{1ad3b4d}, \texttt{82ce12f},
\texttt{1df25f5}) that converts an arbitrary upstream HDL repo into a Harbor task with
correct file-strip manifest, environment Dockerfile, and verifier script. The skill went
through five iterations to handle five different upstream build systems (FuseSoC, Bazel,
Make+Verilator, Icarus, cocotb). We added Ibex, Caliptra, PULP common\_cells, secworks-aes,
VeeR EL2, CVA6, Coral NPU, NVDLA, and OpenPiton. Roughly half of the tasks we tried did not
make the final eleven --- some had upstream build systems that the agent couldn't drive at
all (Caliptra), some had test suites with too much state for the verifier to be
deterministic (CVA6 in some configurations).

\paragraph{The runs.}
The headline run on the eleven tasks took roughly four days of wall-clock and seven
agent$\times$harness combinations: Opus 4.7 and Sonnet 4.6 via Claude Code, GPT-5.5 (effort
\texttt{xhigh}) via Codex CLI, and DeepSeek V4-Pro, GLM-5.1, Kimi K2.6, Qwen 3.6-27B via
OpenCode/OpenRouter. The harness needed several mechanisms we had not anticipated when we
started: per-account concurrency caps because Claude Code rate-limits per subscription, a
disk supervisor because a Sonnet run leaked 812\,GB of Verilator build cache through a
bind-mounted credentials directory and would have crashed the lab server, and a
stuck-job killer because the agents sometimes write at human pace for hours and then
freeze.

\paragraph{The result, briefly.}
The headline numbers are 6/11 perfect for GPT-5.5 and 5/11 for Opus 4.7, with the smaller
open-weight models clustering at 0--2 perfects each. Reading the transcripts we found
that most of the perfect scores came either from sandbox bypasses (CDN mirrors, IP-literal
DNS bypass, vendor-side fetch tools) or from two upstream verifier flaws (an exit-code-only
test in Ibex and an \texttt{assert True} scoreboard in VeeR-EL2). The strong-verifier pass
rate is much lower: 1--2/11 for the best agents, 0/11 for the weakest. The full
discussion is in the companion paper.

\section{What we would do differently}

Three lessons we would carry into the next iteration:

\paragraph{Sandbox first.} We built the harness and added the DNS sinkhole later. By the
time we audited the transcripts, most of the perfect scores had already been bypassed. A
default-deny egress firewall with explicit allowlists, plus an audit pass that flags
external network egress in real time, would have caught everything in this paper before
the run completed.

\paragraph{Verifier hardening should be part of task acceptance.} For every task we adopt
from an upstream project, ``does the agent score 1.000 with a stub'' should be a check we
run before greenlighting the task. Two of the three flagged tasks were exploitable in 17
minutes of agent time; we could have caught both during task setup.

\paragraph{Restoration is the right framing for now.} Of the four prior approaches
(PR-replay, CVDP module synthesis, Spec2GDSBench, restoration), restoration is the only
one that simultaneously (a) gives the agent a real test suite, (b) avoids the
specification problem entirely, (c) scales to project-scale work without us having to
write build infrastructure, and (d) produces a benchmark whose top-line we trust.
\textsc{RTL-Universe} as it stands is not the right benchmark forever, but the
restoration framing has more headroom than what we had before.

\paragraph{Acknowledgements.} We thank the lowRISC, chipsalliance, NVIDIA, Princeton, PULP,
Google Coral, and secworks teams for the upstream projects this benchmark is built on.
Repository contributions over this arc are split roughly evenly between the two authors,
with several small contributions appearing under our other GitHub handles (\texttt{Broyojo}
and \texttt{spyre} respectively, plus the \texttt{SpiritSeal} handle that was used during
the early SWE-Universe-Verilog and RTLBench work).

\bibliography{refs}

\end{document}
